\قسمت{معماری چند‌لایه‌ای در وب پیشین و امروز}

معماری چند‌لایه‌ای برنامه‌های کاربردی وب مدرن یا \متن‌لاتین{Client-Side Rendering}
با معماری چند لایه‌ای برنامه‌های کاربردی وب سنتی یا \متن‌لاتین{Server-Side Rendering}
چه تفاوتی دارند؟ با رسم شکل توضیح دهید.

\نمره{۱}

% پاسخ پیشنهادی است و پیش از استفاده در تصحیح نیاز به بازبینی دارد.
\begin{پاسخ}

در معماری سنتی\پانویس{Server-Side Rendering} لایه‌ی نمایش روی سرور است:
مرورگر تقاضا می‌دهد، سرور \متن‌لاتین{HTML} آماده می‌سازد و مرورگر تنها آن را نمایش می‌دهد.

در معماری امروزی\پانویس{Client-Side Rendering} لایه‌ی نمایش به مرورگر منتقل شده است:
سرور یک رابط برنامه‌نویسی داده (معمولا \متن‌لاتین{JSON}) می‌دهد و ساخت صفحه در مرورگر و با جاوا اسکریپت انجام می‌شود.

پیامدها: در حالت دوم بار پردازش نمایش روی کلاینت است، تقاضاهای بعدی سبک‌ترند (تنها داده)،
اما بارگذاری نخست سنگین‌تر است و برای \متن‌لاتین{SEO} به تمهید نیاز دارد.

\شروع{فقرات}
\فقره رسم شکل لایه‌ها برای هر دو حالت: ۵۰ درصد نمره
\فقره مشخص کردن اینکه نمایش کجا ساخته می‌شود: ۵۰ درصد نمره
\پایان{فقرات}

\end{پاسخ}
